\section{三层访问权限：API、Open Weights、Open Source}

\subsection{三层访问模型}
讲者把访问权限分为三层，分别对应三种研究位置：
\begin{itemize}
    \item API access：像“行为科学家”观察黑箱输入输出；
    \item Open weights：像“神经科学家”可探查内部表示；
    \item Open source：像“系统工程师”可重写整个流水线。
\end{itemize}

\begin{table}[H]
\centering
\begin{tabular}{p{2.4cm}p{4.2cm}p{4.8cm}}
\toprule
\textbf{访问层级} & \textbf{典型能力} & \textbf{主要限制} \\
\midrule
API access & Prompt 设计、工具调用编排、黑箱评测 & 难以验证机制成因，难做结构级改造 \\
Open weights & 表征分析、蒸馏/微调、干预实验 & 仍受既有训练配方与架构约束 \\
Open source & 数据-模型-训练-评测全链路重构 & 对工程与算力要求高，复现成本大 \\
\bottomrule
\end{tabular}
\caption{讲者提出的“访问层级”框架及其研究边界}
\end{table}

\begin{importantbox}{同一模型，不同访问层级会导向不同论文类型}
课程里最有价值的提醒是：研究问题并非完全由兴趣决定，也被“可做什么实验”决定。  
API 时代更容易出现策略/提示层论文；Open Source 时代更可能出现架构、训练目标、数据构造与系统 co-design 论文。
\end{importantbox}

\subsection{为什么要把 Open Weights 与 Open Source 区分开}
\begin{knowledgebox}{Open Weights 不等于 Open Science}
拿到权重可以做很多事情，但你仍不知道：
\begin{itemize}
    \item 数据混合比例如何影响能力边界；
    \item 训练课程（curriculum）和对齐策略如何改变行为；
    \item 性能提升来自架构还是数据工程。
\end{itemize}
因此，Open Weights 是重要中间态，但并未抵达“可完整验证”的科学状态。
\end{knowledgebox}

\begin{warningbox}{政策讨论中的概念混用风险}
产业讨论里经常把“可调用 API”“可下载权重”“可复现训练”混成一个词“开放”。这会让监管与协作机制设计偏离真实技术约束，造成错误激励。
\end{warningbox}

\subsection{本章小结}
讲者提出的三层框架能把“开放性争论”从口号转为工程问题：不同层级对应不同实验能力与不同风险暴露，不能混为一谈。


